% === chapters/01_mechanics/01_kinematics.tex ===
\section{Кинематика}

% --- Задача 1 ---
\problem{1}{800}
Тело брошено вертикально вверх с начальной скоростью $v_0 = 20$ м/с. 
На какую максимальную высоту оно поднимется? Сопротивлением воздуха пренебречь.

\begin{solution}
Используем формулу: $h = \dfrac{v_0^2}{2g} = \dfrac{20^2}{2 \cdot 9.8} \approx 20.4$ м.

\textbf{Ответ:} $20.4$ м.
\end{solution}

% --- Задача 2 ---
\problem{2}{1200}
Два тела движутся навстречу друг другу из точек $A$ и $B$, расстояние между которыми $L = 100$ м. 
Скорости тел $v_1 = 5$ м/с и $v_2 = 3$ м/с. Через какое время они встретятся?

\begin{solution}
Время встречи: $t = \dfrac{L}{v_1 + v_2} = \dfrac{100}{5 + 3} = 12.5$ с.

\textbf{Ответ:} $12.5$ с.
\end{solution}

% --- Задача 3 ---
\problem{3}{1500}
Автомобиль движется равноускоренно из состояния покоя. За четвёртую секунду он прошёл путь $s = 14$ м. 
Найдите ускорение автомобиля.

\begin{solution}
Путь за $n$-ю секунду: $s_n = v_0 + \dfrac{a}{2}(2n-1)$.

Так как $v_0 = 0$, то $s_4 = \dfrac{a}{2}(2 \cdot 4 - 1) = \dfrac{7a}{2}$.

Отсюда $a = \dfrac{2s_4}{7} = \dfrac{2 \cdot 14}{7} = 4$ м/с².

\textbf{Ответ:} $4$ м/с².
\end{solution}